% Part: first-order-logic
% Chapter: models-theories
% Section: size-of-structures

\documentclass[../../../include/open-logic-section]{subfiles}

\begin{document}

\olfileid{fol}{mat}{siz}

\olsection{التعبير عن حجم ال\printtoken{P}{structure}}

\begin{explain}
قد نعبر عن بعض خواص البنى من غير حاجة إلى رمز غير منطقي.
فمن ال!!{sentence}s ما يصدق في !!a{structure} إذا وفقط إذا كان
!!{domain} تلك ال!!{structure} مشتملًا على عدد معيّن~$n$
من ال!!{element}s، إما على الأقل، وإما على الأكثر، وإما بالضبط.
\end{explain}

\begin{prop}
تصدق ال!!{sentence}
\begin{multline*}
  !A_{\ge n} \ident \lexists[x_1][\lexists[x_2][\dots\lexists[x_n][{}]]]\\
  \begin{aligned}
  (\eq/[x_1][x_2] \land {}
  \eq/[x_1][x_3] \land \eq/[x_1][x_4] \land \dots \land \eq/[x_1][x_n] \land {}\\
\eq/[x_2][x_3] \land \eq/[x_2][x_4] \land \dots \land {} \eq/[x_2][x_n] \land {} \\
\vdots\\
\eq/[x_{n-1}][x_n])
\end{aligned}
\end{multline*}
في !!a{structure}~$\Struct M$ إذا وفقط إذا كان $\Domain M$ يحتوي على
$n$ من ال!!{element}s على الأقل. وبناء عليه، يكون
$\Sat{M}{\lnot !A_{\ge n+1}}$ إذا وفقط إذا كان $\Domain M$ يحتوي على
$n$ من ال!!{element}s على الأكثر.

\end{prop}

\begin{prop}
تصدق ال!!{sentence}
\begin{multline*}
!A_{= n} \ident \lexists[x_1][\lexists[x_2][\dots\lexists[x_n][{}]]] \\
\begin{aligned}
  (\eq/[x_1][x_2] \land {}
  \eq/[x_1][x_3] \land \eq/[x_1][x_4] \land \dots \land \eq/[x_1][x_n] \land {}\\
\eq/[x_2][x_3] \land \eq/[x_2][x_4] \land \dots \land {} \eq/[x_2][x_n] \land {} \\
\vdots\\
\eq/[x_{n-1}][x_n] \land {} \\
\lforall[y][(\eq[y][x_1] \lor \dots \lor \eq[y][x_n])]
\end{aligned}
\end{multline*}
في !!a{structure}~$\Struct M$ إذا وفقط إذا كان $\Domain M$ يحتوي على
$n$ من ال!!{element}s بالضبط.
\end{prop}

\begin{prop}
تكون ال!!{structure} لا نهائية إذا وفقط إذا كانت نموذجًا لـ
\[
\{!A_{\ge 1}, !A_{\ge 2}, !A_{\ge 3}, \dots \}.
\]
\end{prop}

وأما جملة منطقية صرفة واحدة تصدق في~$\Struct M$ إذا وفقط إذا
كان $\Domain M$ لا نهائيًا، فلا توجد. ومع ذلك يمكن وضع
!!{sentence}s ذات !!{predicate}s غير منطقية لا يكون لها إلا
نماذج لا نهائية، وإن لم تكن كل !!{structure} لا نهائية نموذجًا لها.
وأما الانتهاء، وكون البنية !!{nonenumerable}، فلا يمكن التعبير
عن أي منهما حتى بمجموعة لا نهائية من ال!!{sentence}s.
وهذه النتائج من مقتضيات مبرهنتي التراص ولوفنهايم--سكولم.

\end{document}
